% from ega0-17-fr.tex
\oldpage[0\textsubscript{IV-1}]{43}
  foncteurs dérivés que si $\operatorname{dim.proj}(M)\leq n$ (resp.
  $\operatorname{dim.inj}(M)\leq n$), on a $L^iT(M)=0$ (resp.
  $R^iT(M)=0$) \emph{pour tout $i>n$}.
  \item[\rm (iii)] Si on suppose $A$ \emph{noethérien} et $M$ \emph{de type
  fini}, la dernière interprétation de la dimension projective donnée dans
  (17.2.1) montre que si $\operatorname{dim.proj}(M)=n$, $M$ admet une
  résolution gauche de longueur $n$ formée de modules projectifs \emph{de type
  fini}. Si de plus $A$ est un anneau \emph{local} noethérien, $M$ admet une
  résolution de longueur $n$ par des modules \emph{libres de type fini}, un
  $A$-module projectif de type fini étant alors libre (Bourbaki,
  \emph{Alg.~comm.}, chap.~II, \S~3, no~2, cor.~2 de la prop.~5).
\end{enumerate}
\end{remarks}

\begin{lemma}[17.2.3]
\label{0.17.2.3-fr}
Soient $A$ un anneau, $M$ un $A$-module. Pour que
$\operatorname{dim.inj}(M)\leq n$, il faut et il suffit que pour tout
$A$-module monogène $N$, on ait $\operatorname{Ext}^{n+1}_A(N,M)=0$.
\end{lemma}

\begin{proof}
Avec les notations de (17.2.1), il s'agit de prouver que $C$ est injectif~;
pour tout $A$-module $N$, $\operatorname{Ext}^{n+1}_A(N,M)$ est isomorphe à
$\operatorname{Ext}^1_A(N,C)$ (M, V, 7), donc on a
$\operatorname{Ext}^1_A(N,C)=0$ pour tout $A$-module $N$ de la forme
$A/\mathfrak{J}$, où $\mathfrak{J}$ est un idéal quelconque de $A$. La suite
exacte des Ext montre alors que l'homomorphisme canonique
$\operatorname{Hom}(A,C)\longrightarrow\operatorname{Hom}(\mathfrak{J},C)$ est
surjectif pour tout idéal $\mathfrak{J}$ de $A$, ce qui implique que $C$ est un
$A$-module injectif (M, I, 3.2).
\end{proof}

\begin{lemma}[17.2.4]
\label{0.17.2.4-fr}
Soient $A$ un anneau noethérien, $M$ un $A$-module de type fini. Pour que
$\operatorname{dim.proj}(M)\leq n$, il faut et il suffit que pour tout
$A$-module monogène $N$, on ait $\operatorname{Ext}^{n+1}_A(M,N)=0$.
\end{lemma}

\begin{proof}
On sait en effet (M, VI, 2.5) que la condition
$\operatorname{dim.proj}(M)\leq n$ est équivalente à
$\operatorname{Ext}^{n+1}_A(M,N)=0$ pour tout $A$-module $N$ \emph{de type
fini}~; pour voir que la condition de l'énoncé est aussi suffisante, on
raisonne par récurrence sur le nombre de générateurs $m$ de $N$~: il y a un
sous-module $N_1$ de $N$ engendré par $m-1$ éléments et tel que
$N_2=N/N_1$ soit monogène~; de la suite exacte
$0\longrightarrow N_1\longrightarrow N\longrightarrow N_2\longrightarrow0$,
on déduit alors la suite exacte
$\operatorname{Ext}^{n+1}_A(M,N_1)\longrightarrow
\operatorname{Ext}^{n+1}_A(M,N)\longrightarrow
\operatorname{Ext}^{n+1}_A(M,N_2)$ et l'hypothèse de récurrence montre que la
condition de l'énoncé entraîne bien $\operatorname{Ext}^{n+1}_A(M,N)=0$.
\end{proof}

\begin{corollary}[17.2.5]
\label{0.17.2.5-fr}
Soient $A$ un anneau noethérien, $M$ un $A$-module. On a
\begin{equation}
\label{0.17.2.5.1-fr}
\tag{17.2.5.1}
  \operatorname{dim.inj}_A(M)
  =\sup_{\mathfrak{m}}
  \bigl(\operatorname{dim.inj}_{A_{\mathfrak{m}}}(M_{\mathfrak{m}})\bigr)
\end{equation}
et si $M$ est \emph{de type fini}
\begin{equation}
\label{0.17.2.5.2-fr}
\tag{17.2.5.2}
  \operatorname{dim.proj}_A(M)
  =\sup_{\mathfrak{m}}
  \bigl(\operatorname{dim.proj}_{A_{\mathfrak{m}}}(M_{\mathfrak{m}})\bigr)
\end{equation}
où $\mathfrak{m}$ parcourt l'ensemble des idéaux premiers (ou l'ensemble des
idéaux maximaux) de $A$.
\end{corollary}

\begin{proof}
En effet, si $N$ est un $A$-module de type fini (donc de présentation finie),
on a, pour toute partie multiplicative $S$ de $A$ et pour tout $i\geq0$
\[
  S^{-1}\operatorname{Ext}^i_A(N,M)
  =\operatorname{Ext}^i_{S^{-1}A}(S^{-1}N,S^{-1}M)
\]
par platitude, en considérant une résolution libre de $M$ et en utilisant le
fait que la relation précédente est vraie pour $i=0$ (Bourbaki,
\emph{Alg.~comm.}, chap.~II, \S~2, no~7, prop.~19). En particulier
$\operatorname{Ext}^i_{A_{\mathfrak{m}}}
(A_{\mathfrak{m}}/\mathfrak{J}A_{\mathfrak{m}},M_{\mathfrak{m}})
=(\operatorname{Ext}^i_A(A/\mathfrak{J},M))_{\mathfrak{m}}$ pour tout idéal
premier $\mathfrak{m}$ et tout idéal $\mathfrak{J}$ de $A$~; si l'on tient
compte de (17.2.3) et de ce que tout
